%-------------------------------------------------------------------------------
%   CHAPTER 4 - EXPERIMENTS AND RESULTS (SHORTENED)
%-------------------------------------------------------------------------------

\section{Туршилт, үр дүн}
\label{sec:experiments}

Тус бүлэгт Бүлэг 1-д дэвшүүлсэн RQ1--RQ3 асуултад хариулах туршилтын үр дүнг танилцуулна: (i) FCFS claim загвар \texttt{FOR UPDATE SKIP LOCKED}-оор race condition-ыг шийдэж буй эсэх (RQ1); (ii) RLS нь multi-tenant isolation-ыг application-layer алдаанаас үл хамааран хангаж буй эсэх (RQ2); (iii) системийн performance (p95 < 2~с, concurrency safe) шаардлагад нийцэж буй эсэх (RQ3-ийн дэд асуулт). Туршилтыг concurrency, load, penetration, functional (E2E) гэсэн дөрөвт хуваав.

%===============================================================================
\subsection{Туршилтын орчин ба аргачлал}
\label{subsec:test_environment}

Туршилтыг production-аас тусгаарласан staging орчинд явуулсан. Орчны үзүүлэлт: Supabase Pro tier, PostgreSQL 15.6, PgBouncer transaction pooling, pool size 15. Клиент машин: MacBook Pro (Apple M1, 16~GB RAM), Node.js 20.11. Клиент — Supabase дундаж RTT 35--45~ms. Production schema, RLS бодлого, RPC функц, триггерүүдийг staging-т хуулбарласан.

Туршилтын өгөгдлийг \texttt{@faker-js/faker}-аар үүсгэв: 20 байгууллага, 200 хэрэглэгч (100 борлуулагч + 100 хүргэгч), 500 бүтээгдэхүүн, 1{,}000 захиалга, 1{,}000 даалгавар. Хүргэгчдийн 80\%-ийг \texttt{approved}, 20\%-ийг \texttt{pending} төлөвт санаатай үлдээсэн нь RLS-ийн approved-only шүүлтийг шалгах зорилготой. Туршилт бүрийн өмнө \texttt{TRUNCATE ... RESTART IDENTITY CASCADE}-ээр өгөгдлийг анхны төлөвт буцаасан.

Ашигласан хэрэгсэл: \texttt{@supabase/supabase-js 2.46} (concurrency test, \texttt{Promise.all()}), \texttt{k6 0.54}~\cite{k6docs} (load test, p50/p95/p99 хэмжилт), \texttt{psql} + \texttt{SET request.jwt.claims} (RLS verification), захиалгат JS script-ууд (penetration, OTP brute-force, E2E). Бүх script, seed файл, k6 тест Git-ийн \texttt{tests/} санд байршиж давтагдах боломжтой.

%===============================================================================
\subsection{Зэрэгцээ хандалтын туршилт}
\label{subsec:concurrency_test}

Зорилго: $N$ хүргэгч нэг даалгаврыг зэрэг claim хийхэд зөвхөн нэг нь амжилттай болж, үлдсэн нь чимээгүйгээр татгалзагдах эсэхийг шалгах. Арга зүй: (1)~шинэ \texttt{delivery\_tasks} мөр \texttt{status='open'}-оор үүсгэх, (2)~$N$ хүргэгчийн JWT токеноор client instance бэлтгэх, (3)~\texttt{Promise.all()}-оор $N$ RPC дуудлагыг зэрэг илгээх, (4)~өгөгдлийн сангаас \texttt{SELECT courier\_id}-г шалгаж нэг мөр л бичигдсэн invariant-ыг батлах. $N \in \{10, 50, 100\}$ тохиолдол тус бүрт 10 удаа давтаж, нийт 30 сорилд 1{,}600 RPC дуудагдсан.

\begin{table}[H]
\centering
\renewcommand{\arraystretch}{1.25}
\caption{Зэрэгцээ claim туршилтын үр дүн ($N$ хүргэгч, 10 давталт)}
\label{tab:concurrency_results}
\footnotesize
\begin{tabular}{|c|c|c|c|c|c|}
\hline
$N$ & Давталт & Амжилтын дундаж & Давхардсан & Амжилтын \% & Дундаж хугацаа \\
\hline
10  & 10 & 1.00 & 0 & 100\% & 124~ms \\
\hline
50  & 10 & 1.00 & 0 & 100\% & 287~ms \\
\hline
100 & 10 & 1.00 & 0 & 100\% & 512~ms \\
\hline
\end{tabular}
\end{table}

Нийт 30 сорилын туршид \textbf{давхардсан claim нэг ч удаа гараагүй}. Үлдсэн $N{-}1$ хүсэлт \texttt{NULL} буцаасан. Клиент тал үүнийг ``даалгавар авагдсан'' мэдэгдэл болгон хөрвүүлнэ. 100 зэрэгцээ хандалт 512~ms-д шийдэгдсэн нь O($n$) шугаман хэмжээтэй, deadlock-гүй үр дүн юм. Энэ нь pessimistic locking-ийн онолын хүлээгдсэн зан төлөвтэй нийцнэ~\cite{bernstein1987concurrency, kleppmann2017ddia}.

Хэрэв claim нь ердийн \texttt{UPDATE ... WHERE status='open'}-оор хэрэгжсэн бол MVCC-ийн улмаас хоёр зэрэгцээ UPDATE хоёулаа мөрийг ``нээлттэй'' гэж үзэж \textit{lost update} үүсэх магадлалтай. Иймд \texttt{FOR UPDATE SKIP LOCKED} нь зөвхөн performance-ийн биш, \textit{correctness}-ийн хувьд зайлшгүй шийдэл юм.

%===============================================================================
\subsection{Гүйцэтгэлийн туршилт (Load Test)}
\label{subsec:load_test}

Гүйцэтгэлийн туршилтыг k6-аар 100 VU, 10~мин ramp-up + 20~мин тогтмол ачаалалтайгаар ажиллуулсан. Үр дүнг Хүснэгт~\ref{tab:endpoint_latency}-д нэгтгэв.

\begin{table}[H]
\centering
\renewcommand{\arraystretch}{1.25}
\caption{k6 ачааллын туршилтын endpoint-үүдийн хариу хугацаа (ms)}
\label{tab:endpoint_latency}
\footnotesize
\begin{tabular}{|p{4.8cm}|c|c|c|c|c|}
\hline
\textbf{Endpoint} & \textbf{p50} & \textbf{p95} & \textbf{p99} & \textbf{Max} & \textbf{Амжилт} \\
\hline
\texttt{GET available\_tasks}      & 145 & 312 & 456 & 892   & 100\% \\
\hline
\texttt{POST claim\_delivery\_task} & 287 & 621 & 834 & 1\,200 & 100\% \\
\hline
\texttt{POST verify\_epod\_otp}    & 198 & 287 & 345 & 523   & 100\% \\
\hline
\texttt{GET orders} (борлуулагч)   & 167 & 389 & 512 & 734   & 100\% \\
\hline
Dashboard KPI агрегат              & 234 & 567 & 789 & 1\,100 & 100\% \\
\hline
\end{tabular}
\end{table}

Хүснэгтийн үр дүнг визуал хэлбэрээр Зураг~\ref{fig:endpoint_latency_chart}-д харуулав.

\begin{figure}[H]
\centering
\resizebox{0.85\textwidth}{!}{%
\begin{tikzpicture}
\begin{axis}[
    ybar,
    bar width=12pt,
    width=0.95\textwidth,
    height=7.5cm,
    enlarge x limits=0.12,
    ymin=0, ymax=900,
    ylabel={Хариу хугацаа (ms)},
    symbolic x coords={available\_tasks,claim\_task,verify\_otp,orders,dashboard},
    xtick=data,
    x tick label style={font=\small, rotate=15, anchor=east},
    ymajorgrids=true,
    grid style={gray!30},
    legend style={at={(0.5,-0.25)}, anchor=north, legend columns=3, font=\small},
    ylabel style={font=\small},
    tick label style={font=\small}
]
\addplot[fill=blue!40, draw=blue!60!black] coordinates {
    (available\_tasks,145) (claim\_task,287) (verify\_otp,198) (orders,167) (dashboard,234)
};
\addplot[fill=orange!60, draw=orange!70!black] coordinates {
    (available\_tasks,312) (claim\_task,621) (verify\_otp,287) (orders,389) (dashboard,567)
};
\addplot[fill=red!50, draw=red!70!black] coordinates {
    (available\_tasks,456) (claim\_task,834) (verify\_otp,345) (orders,512) (dashboard,789)
};
\legend{p50, p95, p99}
\end{axis}
\end{tikzpicture}
}
\caption{k6 load test — 5 endpoint-ийн p50/p95/p99 latency хуваарилалт (100 VU, 30 мин)}
\label{fig:endpoint_latency_chart}
\end{figure}

Гол ажиглалтууд:
\begin{itemize}
  \item \texttt{available\_tasks} нь \texttt{(status)} талбарын partial index ашигладаг тул p95 312~ms-д багтсан.
  \item \texttt{claim\_delivery\_task}-ийн p95 621~ms нь lock acquisition-оос ихээр хамаарч, \S\ref{subsec:concurrency_test}-ийн үр дүнтэй нийцнэ.
  \item \texttt{verify\_epod\_otp}-ийн p50 198~ms-ын ойролцоогоор 160~ms нь bcrypt (cost 10) hash comparison-д зарцуулагдана. Энэ нь brute-force эсэргүүцэхэд зориудаар тавьсан зардал юм~\cite{provos1999bcrypt}.
  \item Бүх endpoint \textbf{p95 < 2~с} (хатуу хязгаар), \textbf{p99 < 1.5~с} (зөөлөн хязгаар)-ыг хангасан. \texttt{claim}-ийн max 1.2~с нь edge case-ийг анхаарах шаардлагатайг илтгэнэ.
\end{itemize}
Иймд 100 VU ачаалал дор performance шаардлага биелсэн гэж дүгнэв.

%===============================================================================
\subsection{Аюулгүй байдлын туршилт}
\label{subsec:security_test}

\subsubsection{RLS бодлогын нэвтрэлтийн тест}

Org~A хэрэглэгчийн JWT-ээр Org~B-гийн өгөгдөлд хандах 6 хувилбарыг туршив. Хүлээгдэж буй үр дүн: бүх тохиолдолд RLS хандалтыг хаах, эсвэл 0 мөр буцаах.

\begin{table}[H]
\centering
\renewcommand{\arraystretch}{1.25}
\caption{RLS бодлогын нэвтрэлтийн тестийн үр дүн}
\label{tab:rls_pentest}
\footnotesize
\begin{tabular}{|c|p{7cm}|p{4cm}|c|}
\hline
\textbf{ID} & \textbf{Хувилбар} & \textbf{Бодит үр дүн} & \textbf{Төлөв} \\
\hline
T1 & Org~A → Org~B \texttt{orders} SELECT & 0 мөр & $\checkmark$ \\
\hline
T2 & Org~A → Org~B \texttt{delivery\_tasks} UPDATE & 0 мөр шинэчлэгдсэн & $\checkmark$ \\
\hline
T3 & \texttt{pending} хүргэгч \texttt{available\_tasks} SELECT & 0 мөр & $\checkmark$ \\
\hline
T4 & Хүргэгч~A → Хүргэгч~B \texttt{courier\_earnings} SELECT & 0 мөр & $\checkmark$ \\
\hline
T5 & SQL injection (\texttt{.eq('org\_id', "' OR 1=1 --")}) & Parameterized, 0 мөр & $\checkmark$ \\
\hline
T6 & Алдаатай / хуурамч JWT & 401 Unauthorized & $\checkmark$ \\
\hline
\end{tabular}
\end{table}

Бүх 6 тест тэнцсэн. T5 нь PostgREST-ийн prepared statement хандлагыг, T3 нь approved шалгалт клиент биш ӨС талд хийгддэгийг тус тус баталсан.

\subsubsection{OTP brute-force ба JWT баталгаажуулалт}

Идэвхтэй OTP сесст буруу кодыг дараалан 6 удаа илгээв. Тав дахь оролдлогын дараа өгөгдлийн сангийн түвшинд блоклогдсон. Зургаа дахь хүсэлт нь bcrypt comparison хүртэл хүрэлгүй 43~ms-д эрт таслагдсан (1--5 оролдлого: 72--86~ms; 6-р оролдлого: ``Хязгаар хэтэрсэн'' алдаа). Brute-force амжилтын онолын магадлал $P = 5/10^6 = 5 \times 10^{-4}\%$. OTP-ийн 10~мин expiry-тэй хослуулахад практикт мэдэгдэхүйц бус~\cite{owasp2024authcheatsheet}. Expiry шалгах нэмэлт тестэд 11~минутын дараа зөв кодоор илгээсэн хүсэлт ``expired'' алдаагаар татгалзагдсан нь \texttt{expires\_at} шалгалтыг баталсан.

JWT-ийн 4 хэлбэрийн дайралтыг шалгав:
\begin{enumerate}
  \item \texttt{org\_id} claim-ийг өөрчилсөн;
  \item \texttt{role='service\_role'} болгож privilege escalation оролдсон;
  \item \texttt{alg=none} signature-гүй токен;
  \item хугацаа дууссан токен.
\end{enumerate}
Бүгд PostgREST түвшинд 401 Unauthorized-аар татгалзагдсан. Өгөгдлийн сан хүртэл хүсэлт дамжаагүй~\cite{jones2015jwt}.

%===============================================================================
\subsection{End-to-End (E2E) туршилт}
\label{subsec:functional_test}

\S\ref{subsec:system_overview}-ийн 8 алхамт үндсэн урсгалыг end-to-end script-ээр туршиж, цагийн тэмдгүүдийн invariant (\texttt{assigned\_at} < \texttt{picked\_up\_at} < \texttt{delivered\_at} < \texttt{completed\_at}) зөрөлгүй биелсэн. Үр дүнг Хүснэгт~\ref{tab:e2e_flow}-д харуулав.

\begin{table}[H]
\centering
\renewcommand{\arraystretch}{1.2}
\caption{E2E захиалгын урсгалын баталгаажуулалт}
\label{tab:e2e_flow}
\footnotesize
\begin{tabular}{|c|p{4.5cm}|p{5.5cm}|c|}
\hline
\textbf{\#} & \textbf{Алхам} & \textbf{Шалгуур} & \textbf{Төлөв} \\
\hline
1 & Байгууллага бүртгэл & \texttt{orgs} мөр + JWT \texttt{org\_id} claim & $\checkmark$ \\
\hline
2 & Бүтээгдэхүүн нэмэх & \texttt{products} + RLS шүүлт & $\checkmark$ \\
\hline
3 & Захиалга хүлээн авах & \texttt{orders.status=paid} & $\checkmark$ \\
\hline
4 & Даалгавар нийтлэх & \texttt{delivery\_tasks.status=open} & $\checkmark$ \\
\hline
5 & Хүргэгч авах & \texttt{status=assigned}, \texttt{assigned\_at} & $\checkmark$ \\
\hline
6 & Бараа авах & \texttt{status=picked\_up}, \texttt{picked\_up\_at} & $\checkmark$ \\
\hline
7 & Хүргэх & \texttt{status=delivered}, \texttt{delivered\_at} & $\checkmark$ \\
\hline
8 & ePOD + орлого & \texttt{status=completed} + \texttt{courier\_earnings} мөр & $\checkmark$ \\
\hline
\end{tabular}
\end{table}

Серверийн бодит тооцооллын хугацаа 2.4~секунд, нийт урсгал (клиент sleep-тэй) 6.8~секунд байлаа. ePOD дэд урсгал нь 5 алхмыг зөв дарааллаар гүйцэтгэв: OTP үүсгэх → bcrypt hash хадгалах → \texttt{verify\_epod\_otp} RPC → completed trigger → \texttt{record\_courier\_earning} trigger. Энэ нь \S\ref{subsec:security_test}-ийн brute-force хамгаалалттай нэгдмэл ажиллаж буйг баталсан.

%===============================================================================
\subsection{Үр дүнгийн дүгнэлт}
\label{subsec:results_summary}

\subsubsection{Хязгаарлалт}

Туршилт нь (i) Supabase Pro tier-ийн нөөцийн хязгаарт, (ii) synthetic өгөгдөл дээр (peak cyclicity гэх мэт бодит зан төлөвгүй), (iii) single-region тохиргоонд, (iv) 100 VU хүртэлх ачаалалд, (v) тогтвортой сүлжээнд (RTT 35--45~ms) хийгдсэн. 500+ VU, multi-region, 3G/4G сүлжээний нөхцөл үнэлэгдээгүй тул эдгээр нь ирээдүйн өргөтгөлийн судалгааны чиглэл болно.

\subsubsection{Шаардлагын биелэлтийн нэгтгэл}

\begin{table}[H]
\centering
\renewcommand{\arraystretch}{1.25}
\caption{Системийн чанарын шаардлагуудын туршилтаар баталгаажсан биелэлт}
\label{tab:requirement_verification}
\footnotesize
\begin{tabular}{|p{2.5cm}|p{5.5cm}|p{4cm}|c|}
\hline
\textbf{Ангилал} & \textbf{Шаардлага} & \textbf{Баталгаа} & \textbf{Төлөв} \\
\hline
Гүйцэтгэл & available\_tasks < 2~с & Хүснэгт~\ref{tab:endpoint_latency} (p95=312~ms) & $\checkmark$ \\
\hline
Гүйцэтгэл & claim < 1~с & Хүснэгт~\ref{tab:endpoint_latency} (p99=834~ms) & $\checkmark$ \\
\hline
Аюулгүй байдал & JWT validation & \S\ref{subsec:security_test} (4 хувилбар) & $\checkmark$ \\
\hline
Аюулгүй байдал & RLS тусгаарлалт & Хүснэгт~\ref{tab:rls_pentest} (T1--T5) & $\checkmark$ \\
\hline
Аюулгүй байдал & bcrypt + HTTPS & OTP тулгалт + TLS 1.3 & $\checkmark$ \\
\hline
Өргөтгөх боломж & Зэрэгцээ claim хамгаалалт & Хүснэгт~\ref{tab:concurrency_results} (100 VU, 0 давхардал) & $\checkmark$ \\
\hline
Өргөтгөх боломж & Олон байгууллага зэрэг & T1, T4 + E2E & $\checkmark$ \\
\hline
Хүртээмж & Сессийн автомат сэргээлт & E2E: токен refresh зөв & $\checkmark$ \\
\hline
Ашиглах чадвар & Хүргэгч < 5 мин бүртгэл & E2E: ~3 мин & $\checkmark$ \\
\hline
\end{tabular}
\end{table}

Системийн чанарын шаардлагууд туршилтаар бүрэн биелж, RQ1--RQ3 гурван асуултад тоон баталгаа олгосон.

\vspace{0.5cm}

%-------------------------------------------------------------------------------
%   SECTION - CONCLUSION
%-------------------------------------------------------------------------------

\section{Дүгнэлт}
\label{sec:conclusion}

Төгсөлтийн ажлын хүрээнд Монголын ЖДБ хоорондын захиалга, хүргэлтийн процессыг цахимжуулах \textbf{``Байгууллагын бүтээгдэхүүн хүргэлтийн нэгдсэн систем''}-ийг боловсруулж, туршилтын үр дүнгээр баталгаажуулав.

%===============================================================================
\subsection{Зорилтуудын хэрэгжилт}

Дэвшүүлсэн дөрвөн зорилтыг дараах байдлаар хэрэгжүүлсэн. Туршилтын үр дүнд үндэслэн зорилт бүрийн гүйцэтгэлд холбогдох нотолгоог нэмж оруулсан.

\begin{table}[H]
\centering
\renewcommand{\arraystretch}{1.3}
\caption{Зорилтуудын хэрэгжилтийн хураангуй}
\label{tab:objectives_summary}
\begin{tabular}{|c|p{3.5cm}|p{7.5cm}|c|}
\hline
\textbf{\#} & \textbf{Зорилт} & \textbf{Хэрэгжүүлэлт ба нотолгоо} & \textbf{Төлөв} \\
\hline
1 & Судалгаа, шинжилгээ & Uber Freight, Papa Logistics судалж, 4 архитектурын зөрүү тодорхойлсон & $\checkmark$ \\
\hline
2 & Архитектур, загварчлал & 11+ хүснэгт бүхий ӨС, 2 клиент (мобайл + веб) + Supabase backend, диаграммууд & $\checkmark$ \\
\hline
3 & Аюулгүй байдал & RLS бодлого, RBAC (4 түвшин), JWT validation, ePOD OTP; RLS penetration 6 тест болон OTP brute-force тестээр бүрэн нотлогдсон (Хүснэгт~\ref{tab:rls_pentest}, \S\ref{subsec:security_test}) & $\checkmark$ \\
\hline
4 & Функционал модулиуд & Захиалга, даалгаврын marketplace, realtime мэдэгдэл, каталог, ePOD, орлого, санхүүгийн тайлан; end-to-end урсгал баталгаажсан (\S\ref{subsec:functional_test}) & $\checkmark$ \\
\hline
\end{tabular}
\end{table}

%===============================================================================
\subsection{Харьцуулсан дүгнэлт}

Судалгаанд авч үзсэн хоёр платформтой харьцуулсан үр дүнг Хүснэгт~\ref{tab:unified_comparison}-д нэгтгэв. Туршилтаар баталгаажсан шинж чанарыг тухайн хүснэгтээр ишилсэн.

\begin{table}[H]
\centering
\renewcommand{\arraystretch}{1.3}
\caption{Гурван системийн нэгдсэн харьцуулалт}
\label{tab:unified_comparison}
\footnotesize
\begin{tabular}{|p{2.8cm}|p{2.8cm}|p{2.8cm}|p{3cm}|c|}
\hline
\textbf{Шалгуур} & \textbf{Uber Freight} & \textbf{Papa Logistics} & \textbf{Миний төсөл} & \textbf{Үнэлгээ} \\
\hline
Даалгавар хуваарилалт & ML алгоритм, recommendation & Оператор + санал & FCFS claim model, concurrency-safe (Хүснэгт~\ref{tab:concurrency_results}) & $\bigstar\bigstar$ \\
\hline
Үнэ тогтоолт & Динамик, ML-д суурилсан & Өрсөлдөөнт bid & Борлуулагч тогтоодог & $\bigstar\bigstar$ \\
\hline
Өгөгдлийн тусгаарлалт & Enterprise түвшин & Application-level & Database-level RLS, 6/6 pentest амжилттай (Хүснэгт~\ref{tab:rls_pentest}) & $\bigstar\bigstar\bigstar\bigstar$ \\
\hline
Audit trail & Execution log & Хязгаарлагдмал & Төлөв шилжилтийн timestamp; бүрэн audit log нь цаашдын ажил & $\bigstar\bigstar$ \\
\hline
API / гадаад integration & Бүрэн, TMS/ERP & Хязгаарлагдмал & RESTful суурь бүтэц & $\bigstar\bigstar\bigstar$ \\
\hline
Realtime хяналт & Location + tracking & Location tracking & Supabase Realtime мэдэгдэл & $\bigstar\bigstar\bigstar$ \\
\hline
Хүргэлтийн баталгаа & Дижитал signature & Тодорхойгүй & ePOD OTP код (brute-force туршигдсан, \S\ref{subsec:security_test}) & $\bigstar\bigstar\bigstar\bigstar$ \\
\hline
Performance & Enterprise түвшин & Үндсэн & p95 < 2с бүх endpoint-д (Хүснэгт~\ref{tab:endpoint_latency}) & $\bigstar\bigstar\bigstar\bigstar$ \\
\hline
Санхүүгийн тайлан & Enterprise түвшин & Үндсэн & Дэлгэрэнгүй (KPI, аналитик) & $\bigstar\bigstar\bigstar\bigstar$ \\
\hline
Өргөтгөх боломж & Cloud, horizontal scaling & Dispatcher bottleneck & Event-driven, horizontal & $\bigstar\bigstar\bigstar\bigstar$ \\
\hline
\end{tabular}
\end{table}

\textbf{Uber Freight-тэй харьцуулбал} чадамжийн ойролцоогоор 35--40\%-д хүрсэн. Uber нь олон жилийн ML, сая+ хэрэглэгчийн өгөгдөл, мянган инженертэй enterprise түвшний систем тул шууд харьцуулах боломжгүй боловч архитектурын суурь зарчим (autonomous claim-based хуваарилалт, database-level security, ePOD verification) ижил чиглэлтэй.

\textbf{Papa Logistics-тэй харьцуулбал} чадамжийн ойролцоогоор 70--75\%-д хүрсэн. Архитектурын хувьд (RLS, audit timestamp, операторгүй marketplace) давуу; харин бодит зах зээлийн туршлага, төлбөрийн integration, location tracking-оор дутмаг.

%===============================================================================
\subsection{Хийж чадаагүй зүйлс ба цаашдын ажил}

Хийгдсэн ажлын дэлгэрэнгүйг §\ref{subsec:results_summary}-ийн Хүснэгт~\ref{tab:objectives_summary} (Зорилтуудын хэрэгжилт) болон Хүснэгт~\ref{tab:requirement_verification} (Шаардлагын биелэлт)-д нэгтгэсэн. Энэ хэсэгт одоогийн хувилбарт хамрагдаагүй, цаашид хөгжүүлэх чиглэлүүдийг тусгайлан дурдав.

\textbf{Захиалагчийн (customer) талын апп.} Одоогийн хувилбарт борлуулагчийн веб апп, хүргэгчийн мобайл апп гэсэн хоёр клиент хөгжүүлсэн. Эцсийн хэрэглэгч (бараа худалдан авагч) нь платформын two-sided marketplace-ийн эрэлтийн тал боловч уг ажлын хүрээнд тусдаа апп хэрэгжээгүй. Одоогоор захиалагч нь борлуулагчийн байгууллагаар дамжуулан захиалга өгдөг хэлбэрээр загварчилсан (\texttt{orders} хүснэгтэд \texttt{customer\_id} талбар бэлтгэгдсэн). Цаашид захиалагчийн мобайл апп (бүтээгдэхүүн хайх, сагс, төлбөр, хүргэлт хянах) нэмэх нь зайлшгүй.

\textbf{Цаашид хөгжүүлэх бусад чиглэл:}
\begin{itemize}
    \item Бодит төлбөрийн integration (QPay, SocialPay, банкны API)
    \item Realtime GPS location tracking
    \item ML-д суурилсан даалгаврын recommendation
    \item Route optimization алгоритм
    \item Push notification систем
    \item Бодит хэрэглэгчтэй pilot туршилт
    \item 100+ concurrent user-тэй өргөтгөсөн load test (read replica, per-tenant sharding нөхцөлд)
\end{itemize}

%===============================================================================
\subsection{Төгсгөл}

Тус төсөл нь Монгол Улсын B2B хүргэлтийн зах зээлд operator-гүй, autonomous FCFS claim хуваарилалт болон database-level security-д суурилсан шинэ архитектурын шийдлийг concurrency, penetration, load туршилтаар тоон үзүүлэлтээр баталгаажуулсан. Цаашид §\ref{subsec:results_summary}-д дурдсан чадамжуудыг нэмж хөгжүүлснээр production-д нэвтрэх боломжтой систем болно.

